\setupcolors[state=start]
\setupcolor[hex]
\definecolor[darkblue][h=003366]
\definecolor[darkgray][h=e6e6e6]
\definecolor[redline][h=d9534f]
\definecolor[blueline][h=5bc0de]
\definecolor[greenline][h=5cb85c]

\setuphead[section][style=\bfa, color=darkblue, before=\blank, after=\blank]
\setuphead[subsection][style=\bf, color=black, before=\blank, after=\blank]
\setuppapersize[A4][A4]
\setupbodyfont[11pt]
\setuplayout[width=middle, height=middle, topspace=2cm, bottomspace=2cm]

\starttext

\title{\bf Comparación de OPP\_v4 y OPP\_v5: pruebas normales equivalentes}
\blank[big]

\startsection[title={Objetivo y alcance}]
Este documento compara las ejecuciones normales de OPP\_v4 y OPP\_v5 que poseen una correspondencia uno a uno. El objetivo es verificar la regresión funcional de la versión v5 y describir las diferencias de tiempo observadas bajo el mismo conjunto de escenarios heredados.

La comparación utiliza exclusivamente las salidas:
\startitemize
\item \type{tests/outputs/opp_v4/normal/output_opp_v4_normal_tests_20260825_192401.csv}.
\item \type{tests/outputs/opp_v5/normal/output_opp_v5_normal_tests_20260825_191715.csv}.
\stopitemize

No se incorporan aún resultados de estrés. Estos deberán añadirse en una sección posterior cuando estén disponibles ambos archivos de salida.
\stopsection

\startsection[title={Criterio de equivalencia}]
Los archivos \type{test_opp_v5_model_01.py} a \type{test_opp_v5_model_19.py}, incluidas las variantes \type{02_01}, \type{02_02} y \type{02_03}, migran uno a uno los escenarios y las aserciones normales de OPP\_v4. En la migración, v5 ejecuta primero el modelo de ruteo y luego el de asignación y secuenciación.

Para aislar el comportamiento heredado de v4, las capacidades nuevas de v5 se neutralizan en esos casos: todos los pickers usan \(V_p=1\) y el límite de frío se fija en un valor no restrictivo. Así, las pruebas comparan factibilidad, rutas, asignación, secuenciación y penalizaciones sobre escenarios equivalentes. Las pruebas v5 específicas de velocidades heterogéneas y cadena de frío (modelos 20 y 21) no se incluyen en la comparación uno a uno.

El tiempo informado es \type{total_time}: el tiempo de ejecución de cada prueba. En v5 incluye ambas etapas del pipeline. No se comparan \type{obj_val} ni \type{model_time} como métricas homogéneas, porque v5 registra como último resultado el de la segunda etapa y su función objetivo no es directamente equivalente a la formulación monolítica de v4.
\stopsection

\startsection[title={Cobertura y éxito de ejecución}]
Se identificaron 24 pares de ejecuciones equivalentes. Veintidós pares finalizaron con estado \type{ok} en ambos modelos; estos son los casos válidos para las métricas temporales. OPP\_v4 registró dos fallos en los casos 10 y 11, mientras que sus pares de v5 finalizaron correctamente.

\blank[medium]
\bTABLE[split=yes, align=middle, frame=on, width=fit]
\setupTABLE[row][1][background=color, backgroundcolor=darkgray, style=bold]
\bTR \bTH Métrica \eTH \bTH OPP\_v4 \eTH \bTH OPP\_v5 \eTH \eTR
\bTR \bTD Ejecuciones normales del CSV \eTD \bTD 24 \eTD \bTD 28 \eTD \eTR
\bTR \bTD Pares equivalentes \eTD \bTD 24 \eTD \bTD 24 \eTD \eTR
\bTR \bTD Pares exitosos en ambos modelos \eTD \bTD 22 \eTD \bTD 22 \eTD \eTR
\bTR \bTD Éxito en los pares equivalentes \eTD \bTD 22 / 24 (91,7\%) \eTD \bTD 24 / 24 (100,0\%) \eTD \eTR
\bTR \bTD Casos adicionales no equivalentes \eTD \bTD 0 \eTD \bTD 4 \eTD \eTR
\eTABLE

Los cuatro casos adicionales de v5 corresponden a pruebas propias de la versión: velocidad diferencial de pickers y restricción de cadena de frío. Se excluyen para no alterar la base comparable.
\stopsection

\startsection[title={Resultados temporales de los pares exitosos}]
Aplicando el criterio de éxito definido en la especificación de análisis, se consideran únicamente las 22 observaciones con \type{status = ok} en ambos modelos. La mejora del promedio está dominada por el escenario 17, que en v4 tardó 6,831 s y en v5 0,310 s. La mediana ofrece una referencia más representativa para los escenarios pequeños y medianos.

\blank[medium]
\bTABLE[split=yes, align=middle, frame=on, width=fit]
\setupTABLE[row][1][background=color, backgroundcolor=darkgray, style=bold]
\bTR \bTH Métrica (s) \eTH \bTH OPP\_v4 \eTH \bTH OPP\_v5 \eTH \bTH Diferencia v5--v4 \eTH \eTR
\bTR \bTD Media \eTD \bTD 0,480 \eTD \bTD 0,214 \eTD \bTD -0,266 (-55,4\%) \eTD \eTR
\bTR \bTD Desvío estándar \eTD \bTD 1,419 \eTD \bTD 0,067 \eTD \bTD -- \eTD \eTR
\bTR \bTD Mediana \eTD \bTD 0,168 \eTD \bTD 0,187 \eTD \bTD +0,019 (+11,1\%) \eTD \eTR
\bTR \bTD Mínimo \eTD \bTD 0,147 \eTD \bTD 0,154 \eTD \bTD -- \eTD \eTR
\bTR \bTD Máximo \eTD \bTD 6,831 \eTD \bTD 0,444 \eTD \bTD -- \eTD \eTR
\eTABLE

En 5 de los 22 pares v5 fue más rápida. En los otros 17, v5 agregó entre 0,003 s y 0,276 s, consistente con el costo fijo de compilar y resolver dos etapas. Por ello, la disminución del promedio no debe interpretarse como una mejora uniforme para las instancias pequeñas: el promedio refleja principalmente el cambio en la instancia de grilla del caso 17.
\stopsection

\startsection[title={Detalle por caso equivalente exitoso}]
La siguiente tabla presenta los 22 pares utilizados en las métricas. El identificador corresponde al sufijo común \type{model\_<id>} de los archivos v4 y v5.

\blank[medium]
\bTABLE[split=yes, align=middle, frame=on, width=broad]
\setupTABLE[row][1][background=color, backgroundcolor=darkgray, style=bold]
\bTR \bTH Caso \eTH \bTH Prueba \eTH \bTH v4 (s) \eTH \bTH v5 (s) \eTH \bTH Δ (s) \eTH \bTH Δ relativo \eTH \eTR
\bTR \bTD 01 \eTD \bTD Capacidad básica \eTD \bTD 0,180 \eTD \bTD 0,271 \eTD \bTD +0,091 \eTD \bTD +50,5\% \eTD \eTR
\bTR \bTD 02 \eTD \bTD Dos pasillos \eTD \bTD 0,161 \eTD \bTD 0,196 \eTD \bTD +0,035 \eTD \bTD +21,5\% \eTD \eTR
\bTR \bTD 02\_01 \eTD \bTD Dos pasillos \eTD \bTD 0,174 \eTD \bTD 0,180 \eTD \bTD +0,006 \eTD \bTD +3,4\% \eTD \eTR
\bTR \bTD 02\_02 \eTD \bTD Un punto por orden \eTD \bTD 0,168 \eTD \bTD 0,444 \eTD \bTD +0,276 \eTD \bTD +164,1\% \eTD \eTR
\bTR \bTD 02\_03 \eTD \bTD Un pasillo \eTD \bTD 0,265 \eTD \bTD 0,186 \eTD \bTD -0,079 \eTD \bTD -29,9\% \eTD \eTR
\bTR \bTD 03 \eTD \bTD Rutas alternativas \eTD \bTD 0,170 \eTD \bTD 0,211 \eTD \bTD +0,041 \eTD \bTD +24,2\% \eTD \eTR
\bTR \bTD 04 \eTD \bTD Ubicaciones repetidas \eTD \bTD 0,164 \eTD \bTD 0,173 \eTD \bTD +0,009 \eTD \bTD +5,5\% \eTD \eTR
\bTR \bTD 05 \eTD \bTD Repetición alternativa \eTD \bTD 0,153 \eTD \bTD 0,191 \eTD \bTD +0,038 \eTD \bTD +24,8\% \eTD \eTR
\bTR \bTD 06 \eTD \bTD Dos ubicaciones \eTD \bTD 0,151 \eTD \bTD 0,178 \eTD \bTD +0,027 \eTD \bTD +17,9\% \eTD \eTR
\bTR \bTD 07 \eTD \bTD Dos misiones \eTD \bTD 0,149 \eTD \bTD 0,179 \eTD \bTD +0,030 \eTD \bTD +19,9\% \eTD \eTR
\bTR \bTD 08 \eTD \bTD Cuatro órdenes \eTD \bTD 0,184 \eTD \bTD 0,183 \eTD \bTD -0,002 \eTD \bTD -0,8\% \eTD \eTR
\bTR \bTD 12 \eTD \bTD Tres órdenes, un picker \eTD \bTD 0,174 \eTD \bTD 0,209 \eTD \bTD +0,034 \eTD \bTD +19,7\% \eTD \eTR
\bTR \bTD 13 \eTD \bTD Seis órdenes, dos pickers \eTD \bTD 0,271 \eTD \bTD 0,261 \eTD \bTD -0,010 \eTD \bTD -3,7\% \eTD \eTR
\bTR \bTD 14 \eTD \bTD Ruta ramificada \eTD \bTD 0,169 \eTD \bTD 0,171 \eTD \bTD +0,003 \eTD \bTD +1,7\% \eTD \eTR
\bTR \bTD 15 \eTD \bTD Ruta con desvío \eTD \bTD 0,155 \eTD \bTD 0,183 \eTD \bTD +0,028 \eTD \bTD +17,9\% \eTD \eTR
\bTR \bTD 16a \eTD \bTD Ruta extendida \eTD \bTD 0,147 \eTD \bTD 0,154 \eTD \bTD +0,006 \eTD \bTD +4,3\% \eTD \eTR
\bTR \bTD 16b \eTD \bTD Penalizaciones \eTD \bTD 0,215 \eTD \bTD 0,300 \eTD \bTD +0,085 \eTD \bTD +39,6\% \eTD \eTR
\bTR \bTD 16c \eTD \bTD Ubicación duplicada \eTD \bTD 0,153 \eTD \bTD 0,179 \eTD \bTD +0,026 \eTD \bTD +17,1\% \eTD \eTR
\bTR \bTD 17 \eTD \bTD Grilla multiubicación \eTD \bTD 6,831 \eTD \bTD 0,310 \eTD \bTD -6,521 \eTD \bTD -95,5\% \eTD \eTR
\bTR \bTD 18 \eTD \bTD Puente desconectado \eTD \bTD 0,192 \eTD \bTD 0,167 \eTD \bTD -0,025 \eTD \bTD -12,9\% \eTD \eTR
\bTR \bTD 19a \eTD \bTD Adelanto \eTD \bTD 0,158 \eTD \bTD 0,191 \eTD \bTD +0,033 \eTD \bTD +20,9\% \eTD \eTR
\bTR \bTD 19b \eTD \bTD Tardanza \eTD \bTD 0,164 \eTD \bTD 0,188 \eTD \bTD +0,025 \eTD \bTD +15,0\% \eTD \eTR
\eTABLE
\stopsection

\startsection[title={Pares excluidos de la estadística temporal}]
Los casos 10 y 11 son equivalentes por diseño, pero no satisfacen el filtro de éxito para v4. Se muestran para mantener trazabilidad, sin usarlos en los promedios ni en el análisis de dispersión.

\blank[medium]
\bTABLE[split=yes, align=middle, frame=on, width=fit]
\setupTABLE[row][1][background=color, backgroundcolor=darkgray, style=bold]
\bTR \bTH Caso \eTH \bTH Estado v4 \eTH \bTH Tiempo v4 (s) \eTH \bTH Estado v5 \eTH \bTH Tiempo v5 (s) \eTH \eTR
\bTR \bTD 10: dos períodos \eTD \bTD failed \eTD \bTD 8,012 \eTD \bTD ok \eTD \bTD 0,750 \eTD \eTR
\bTR \bTD 11: un período \eTD \bTD failed \eTD \bTD 8,018 \eTD \bTD ok \eTD \bTD 0,829 \eTD \eTR
\eTABLE

El estado \type{failed} no implica necesariamente infactibilidad del modelo: identifica una prueba que no completó satisfactoriamente según su aserción. Se requiere revisar sus trazas si se desea atribuir la diferencia a una causa de formulación, implementación o validación.
\stopsection

\startsection[title={Conclusiones y próximos pasos}]
\startitemize[n]
\item {\bf Regresión funcional:} v5 completó los 24 casos heredados equivalentes, incluidos los dos que v4 marcó como fallidos. Los 22 pares exitosos conservan las aserciones de rutas, asignación y calendario heredadas.
\item {\bf Rendimiento normal:} para instancias pequeñas, v5 presenta en general un sobrecosto moderado atribuible al pipeline de dos etapas. La mediana es 0,187 s en v5 frente a 0,168 s en v4.
\item {\bf Instancias estructuradas:} el caso 17 reduce el tiempo de 6,831 s a 0,310 s. Esta observación explica la mejora de 55,4\% en la media, pero debe confirmarse con más instancias antes de generalizarla.
\item {\bf Alcance pendiente:} incorporar los CSV de estrés, separando los escenarios exitosos y reportando su tasa de éxito, tiempos y escalabilidad. Allí será especialmente relevante medir el efecto de descomponer ruteo y scheduling.
\stopitemize
\stopsection

\stoptext
